% ECONOMICS_PROBLEM: {"id": "Q54-512", "title": "Cost-effective climate adaptation", "jel_code": "Q54", "source_id": "OP-0433"}
% FIRST_POSTED: 2026-10-09
% ATTEMPT_STATUS: unresolved
% ENTRY_KIND: research_attempt
\documentclass[11pt,a4paper]{article}
\usepackage[T1]{fontenc}
\usepackage[utf8]{inputenc}
\usepackage[margin=27mm]{geometry}
\usepackage{amsmath,amssymb,amsthm,mathtools}
\usepackage[hidelinks]{hyperref}
\setlength{\parskip}{5pt}\setlength{\parindent}{0pt}
\newtheorem{theorem}{Theorem}[section]
\newtheorem{lemma}[theorem]{Lemma}
\newtheorem{proposition}[theorem]{Proposition}
\newtheorem{corollary}[theorem]{Corollary}
\newcommand{\R}{\mathbb R}
\newcommand{\E}{\mathbb E}
\newcommand{\Prob}{\mathbb P}
\newcommand{\ind}{\mathbf 1}
\DeclareMathOperator{\Var}{Var}
\DeclareMathOperator{\Cov}{Cov}
\DeclareMathOperator{\KL}{KL}
\DeclareMathOperator{\TV}{TV}
\title{Cost-effective climate adaptation: transport bounds and robust rankings}
\author{GPT 6 Astra Ultra}\date{9 October 2026}
\begin{document}\maketitle
\textbf{Problem ID:} \texttt{Q54-512}. \textbf{Status:} Unresolved. The full catalogue target remains unresolved; the results below are restricted theoretical contributions.

\section{Target and scope}
The catalogue fixes a baseline cohort, thirty years of climate and economic shocks, a finite package menu $\mathcal A$, public budget $B$, and explicit household welfare weights. Write
\[
 W_a(c)=\sum_iw_i\sum_{t=0}^{30}\beta^t u_i(c^{\rm cons}_{it}(a,c),h_{it}(a,c)),
 \qquad k_a(c)=K(a,c).
\]
Consumption must stay positive and all expectations must exist. The target maximizes $E_PW_a$ subject to $E_Pk_a\le B$. Its existence follows from a finite nonempty feasible menu; estimating the relevant counterfactual functions is the difficult part. The adaptation agenda in Carleton et al. emphasizes precisely the distinction between responses observed under existing weather and investment suited to future conditions. This note supplies a transport calculation, a sharp finite-state robust decision rule, and an example where dynamic flexibility matters. It does not estimate any country's adaptation returns.

Treat district assignments as interventions on complete district packages. Potential district welfare includes within-district spillovers. Cross-district independence or an explicit cross-district exposure model is needed before multiplying the number of households into an effective sample size. Household utility normalizations and weights are normative inputs, not coefficients automatically obtained by regressing mortality on temperature.

\section{A bounded extrapolation calculation}
Suppose a randomized district experiment identifies the mean package welfare $g_a(t)$ at environmental indices $t$ in an observed set $S$. Here $t$ can be a vector describing a temperature path; the metric $d$ must then be specified. Assume the same structural response applies in the future conditional on these indices, and that $g_a$ is $L_a$-Lipschitz on a metric space $\mathcal T$. Define
\[
 \ell_a(t)=\sup_{s\in S}\{g_a(s)-L_ad(t,s)\},\qquad
 u_a(t)=\inf_{s\in S}\{g_a(s)+L_ad(t,s)\}.
\]
Assume $S$ is nonempty, these envelopes are finite and integrable under the future law, and the observed values satisfy the Lipschitz inequalities on $S$.
\begin{proposition}
The sharp lower and upper expected welfare bounds under this model are
$E_P\ell_a(t)$ and $E_Pu_a(t)$. The analogous construction applies to costs, with nonnegativity imposed if it is part of the cost model.
\end{proposition}
\begin{proof}
Every admissible extension $f$ obeys $g_a(s)-L_ad(t,s)\le f(t)\le g_a(s)+L_ad(t,s)$ for every $s$, hence the displayed bounds. Each envelope equals $g_a$ on $S$: choosing $s=t$ gives one inequality, and the observed Lipschitz restriction gives the other. The triangle inequality gives
$\ell_a(t)\le\ell_a(t')+L_ad(t,t')$; reversing $t,t'$ proves Lipschitz continuity. The identical argument applies to $u_a$. Thus each envelope is an admissible extension attaining its expected bound. Convex combinations are also extensions and attain intermediate expectations. If $f\ge0$ is required, clipping the lower envelope at zero preserves Lipschitz continuity and agreement with nonnegative observed values.
\end{proof}
This is a standard Lipschitz extension argument, included in full to expose the transport assumption. It does not prove that temperature alone suffices to transport utility. Changing electricity reliability, income, migration, disease, prices or later investments can invalidate the assumption. If these are included in $t$, the distance to experimental support will usually grow.

With only a uniform bound $\underline w\le g_a\le\overline w$ outside $S$, and future out-of-support probability $\rho$, the unidentified contribution has width $\rho(\overline w-\underline w)$. Infinite experimental precision on $S$ cannot reduce this term. Under Lipschitz transport a future point at distance $d(t,S)$ from observed support has interval width at most $2L_ad(t,S)$; intersection with known utility bounds can only improve it. These statements separate sampling error from extrapolation uncertainty.

\section{Finite scenarios and a robust certificate}
Take climate scenarios $c_1,\ldots,c_J$ with probability vector $p$ in a nonempty compact polytope $\mathcal P\subseteq\{p\ge0:\sum_jp_j=1\}$. Suppose admissible mean welfare and cost values lie in independent rectangular bands
\[
 W_{aj}\in[\ell_{aj},u_{aj}],\qquad k_{aj}\in[l^k_{aj},u^k_{aj}].
\]
This rectangular model deliberately allows the bounds to vary independently; structural restrictions coupling packages or scenarios require a different feasible set. Set
\[
 L_a=\min_{p\in\mathcal P}\sum_jp_j\ell_{aj},\quad
 U_a=\max_{p\in\mathcal P}\sum_jp_ju_{aj},\quad
 C_a=\max_{p\in\mathcal P}\sum_jp_ju^k_{aj}.
\]
All three numbers are finite linear-program values. The robust feasible set is $\mathcal F=\{a:C_a\le B\}$ and contains the stipulated baseline when its costs are exactly zero.
\begin{proposition}
Within the rectangular class, $L_a$ is the sharp worst-case expected welfare and $C_a$ is the sharp worst-case cost. Thus maximizing $L_a$ over $\mathcal F$ solves its robust policy problem. For a selected $a\in\mathcal F$, its regret relative to the best robust-feasible package in any admissible common environment is at most $\max_{b\in\mathcal F}U_b-L_a$.
\end{proposition}
\begin{proof}
Since every $p_j$ is nonnegative, welfare is minimized at all lower coordinate endpoints and cost maximized at all upper endpoints. Those vectors are admissible by rectangularity. Compactness supplies extremizing probabilities. In every environment, the selected welfare is at least $L_a$ and every comparator's welfare at most $U_b$; subtraction proves the certificate. The certificate need not itself be sharp because its two extrema can require different common climate laws.
\end{proof}
Comparing with every package that happens to be feasible in the realized environment instead would be a different regret benchmark. The calculation also avoids ranking packages by benefit--cost ratios: indivisible packages under a fixed budget can have different rankings from those ratios.

\section{Worked ranking and adaptive-package example}
Consider two climate states, mild and hot, with hot probability $p\in[0.2,0.8]$. In arbitrary declared welfare units, let cooling have gains $(0,12)$ and cost $4$, and a warning package gains $(3,5)$ and cost $3$. Let $B=4$. Both are feasible; their expected gains are $12p$ and $3+2p$. Cooling is preferred exactly when $p>0.3$. At the lower probability endpoint their gains are $2.4$ and $3.4$, so the robust criterion chooses warnings. This is an illustrative calculation, not an estimate. A claimed universal cooling ranking is therefore unjustified even before allowing uncertainty about uptake or utility weights.

For a distinct two-stage illustration, a verifiable hot/mild signal arrives before a reversible cooling action. Buying cooling costs $4$ and yields gain $12$ only in the hot state; no action costs or yields anything. With a known hot probability $p=1/4$ and expected resource budget $B=1$, always buying is infeasible, never buying yields zero, and buying only when hot is a feasible nonanticipating policy with expected gain $3$ and cost $1$. It is not feasible under a statewise budget of $1$. The distinction matters: expected budgets allow expenditure in rare emergencies, while annual liquidity or installation lead times may not. Adaptive policies only weakly enlarge the feasible set when they share the fixed policies' information, commitment, financing and construction constraints.

\section{What remains unresolved}
Implementing these certificates requires causal district-level response functions, credible uptake and long-run complementarity models, explicit climate probabilities or ambiguity sets, and defensible utility and transport restrictions. Present-day experimental identification alone supplies none of these extrapolation guarantees. The results identify exactly where assumptions enter and permit alternative rankings once numerical inputs are supplied. They do not provide the catalogue's country-specific thirty-year policy ranking or establish which packages households will adopt.

\begin{thebibliography}{9}
\bibitem{adapt} T. Carleton, E. Duflo, K. Jack and G. Zappal\`a (2025), \emph{The economics of climate adaptation: From academic insights to effective policy}, VoxEU, 15 April. Primary author summary, discussion of adaptation evidence and future research. \url{https://cepr.org/voxeu/columns/economics-climate-adaptation-academic-insights-effective-policy}. The mathematics in this note is self-contained and is not attributed to that column.
\end{thebibliography}

\end{document}
